\chapter[Band 46: Frankls Vermutung]{Band 46 auf einen Blick}
\noindent\textbf{Frankls Vermutung}\par\medskip
Die allgemeine Frankl-Vermutung wird in diesem Band
\emph{nicht bewiesen}. Bewiesen werden Spezialfälle, Reduktionen und die
Äquivalenz von Mengen- und Halbverbandsform. Vorausgesetzt werden
Mengenfamilien, Quotienten, endliche Kardinalitäten und Band~45.
Die Häufigkeitsaussagen werden wie im Original durch Injektionen formuliert.

\begin{BandUebersicht}
\textbf{Frankl-Familie}
\UebFormel{\begin{aligned}
&\FranklFam M\ \leftrightarrow\\
&M\ne\varnothing\land\bigcup M\ne\varnothing\\
&\quad\land\Finite M\\
&\quad\land\UnionClosedFam(M).\\
&\UnionClosedFam(M)\ \leftrightarrow\\
&\quad\forall X,Y\in M\\
&\qquad(X\cup Y\in M).\\
&M_a=\{X\in M\mid a\in X\},\\
&M_{\bar a}=\{X\in M\mid a\notin X\}.
\end{aligned}}
\UebQuellen{\FormulaRefAuto{Frankl-Familie}[def];
\FormulaRefAuto{Vereinigungsabgeschlossene Familie}[def];
\FormulaRefAuto{Enthaltende Teilfamilie}[def];
\FormulaRefAuto{Vermeidende Teilfamilie}[def]}
&
\textbf{Zerlegung und Endlichkeit der Teilfamilien}
\UebFormel{\begin{aligned}
&M\setminus M_{\bar a}=M_a,\\
&\Finite M\ \vdash\Finite{M_a}\land\Finite{M_{\bar a}}.
\end{aligned}}
\UebQuellen{\FormulaRefAuto{FamContainingAsAvoidingRelativeComplement}[thm];
\FormulaRefAuto{FiniteContainingSubfamily}[thm];
\FormulaRefAuto{FiniteAvoidingSubfamily}[thm]} \\
\addlinespace[.8ex]

\textbf{Halbhäufigkeit und Frankl-Eigenschaft}
Mit den eben definierten Teilfamilien:\UebFormel{\begin{aligned}
&\HalfOcc{a}{M}\ \leftrightarrow\\
&a\in\bigcup M\land
 \exists f:M_{\bar a}\inj M_a,\\
&\Frankl{M}\ \leftrightarrow\exists a\ \HalfOcc{a}{M}.
\end{aligned}}Bei endlichem \(M\) heißt dies: \(a\) liegt in mindestens der
Hälfte der Familienglieder.
\UebQuellen{\FormulaRefAuto{Halbhäufiges Element}[def];
\FormulaRefAuto{Frankl-Eigenschaft}[def]}
&
\textbf{Bewiesener Spezialfall: Zweiermenge}
\UebFormel{\begin{aligned}
&\FranklFam M,\quad a\ne b,\quad \{a,b\}\in M\\
&\qquad\vdash\Frankl{M}.
\end{aligned}}Die bloße Vereinigungsgeschlossenheit einer beliebigen
Frankl-Familie wird dadurch nicht erledigt.
\UebQuellen{\FormulaRefAuto{FranklFamPairMemberFrankl}[thm]} \\
\addlinespace[.8ex]

\textbf{Ununterscheidbarkeitsquotient}
Für \(a,b\in\bigcup M\) setzt man\UebFormel{\begin{aligned}
&a\sim_M b\ \leftrightarrow\\
&\qquad\forall X\in M\ (a\in X\leftrightarrow b\in X).
\end{aligned}}Ersetzt man jedes \(a\in X\) durch seine Klasse \([a]\), entsteht
die Quotientenfamilie \(\OccQuotFam M\).
\UebQuellen{\FormulaRefAuto{SameOccRelDef}[def];
\FormulaRefAuto{OccQuotMapChar}[thm]}
&
\textbf{Reduktion auf endliche Glieder}
\UebFormel{\begin{aligned}
&\FranklFam M\ \vdash\\
&\qquad\FranklFam{\OccQuotFam M},\\
&\Finite M,\ Y\in\OccQuotFam M\\
&\qquad\vdash\Finite Y,\\
&\Frankl{\OccQuotFam M}\ \vdash\Frankl{M}.
\end{aligned}}Es genügt daher, die Vermutung für Frankl-Familien mit
ausschließlich endlichen Gliedern zu beweisen.
\UebQuellen{\FormulaRefAuto{OccQuotFamFranklFam}[thm];
\FormulaRefAuto{OccQuotMembersFinite}[thm];
\FormulaRefAuto{OccQuotFranklTransport}[thm];
\FormulaRefAuto{FranklFiniteMembersReduction}[thm]} \\
\addlinespace[.8ex]

\textbf{Adjunktion der leeren Menge}
\UebFormel{\begin{aligned}
&M^\circ=M\cup\{\varnothing\}.
\end{aligned}}
\UebQuellen{\FormulaRefAuto{FranklEmptyAdjunction}[def]}
&
\textbf{Normalisierung und Rücktransport}
Für \(\FranklFam M\) gilt:
\UebFormel{\begin{aligned}
&\FranklFam M\ \vdash\FranklFam{M^\circ},\\
&\HalfOcc{a}{M^\circ}\ \vdash\HalfOcc{a}{M},\\
&\Frankl{M^\circ}\ \vdash\Frankl{M}.
\end{aligned}}Die globale Vermutung ist äquivalent zu ihrer Einschränkung
auf Familien, welche die leere Menge enthalten.
\UebQuellen{\FormulaRefAuto{FranklEmptyAdjunctionPreservesFamily}[thm];
\FormulaRefAuto{FranklEmptyAdjunctionHalfOccurrenceBack}[thm];
\FormulaRefAuto{FranklEmptyAdjunctionFranklBack}[thm];
\FormulaRefAuto{FranklEmptyFamilyGlobalEquivalence}[thm]} \\
\addlinespace[.8ex]

\textbf{Halbverbandsform}
Für \(\ZeroJoinSemi{A,\JoinOp,0}\) und dessen Ordnung \(\le\)
sei \(\uparrow j=\{x\in A\mid j\le x\}\). Dann bedeutet\UebFormel{\begin{aligned}
&\FranklSemi{A,\JoinOp,0}\ \leftrightarrow\\
&\exists j\in\JoinIrrSet{A,\JoinOp,0}\
 \exists f:\uparrow j\inj(A\setminus\uparrow j).
\end{aligned}}Im endlichen Fall soll also der Hauptfilter eines
Irreduziblen höchstens die halbe Größe von \(A\) haben.
\UebQuellen{\FormulaRefAuto{FranklSemilatticeProperty}[def]}
&
\textbf{Bewiesene Halbverbandsfälle}
Ein \(\JoinOp\)-irreduzibles größtes Element \(t\) genügt:\UebFormel{\begin{aligned}
&\JoinIrr{A,\JoinOp,0}t,\quad
 \forall x\in A\ (x\le t)\\
&\qquad\vdash\FranklSemi{A,\JoinOp,0}.
\end{aligned}}Auch jede endliche, total geordnete solche Struktur mit
größtem Element \(t\ne0\) erfüllt die Eigenschaft.
\UebQuellen{\FormulaRefAuto{JoinIrreducibleTopFrankl}[thm];
\FormulaRefAuto{FranklJoinChains}[thm]} \\
\addlinespace[.8ex]

\textbf{Kanonische Profilfamilie}
Für einen endlichen Halbverband \((A,\JoinOp,0)\) mit Null setze
\(J=\JoinIrrSet{A,\JoinOp,0}\) und\UebFormel{\begin{aligned}
&J(x)=\{j\in J\mid j\le x\},\\
&\mathcal U=\{J\setminus J(x)\mid x\in A\},\\
&\Phi(x)=J\setminus J(x).
\end{aligned}}
\UebQuellen{\FormulaRefAuto{IrreduciblesBelow}[def];
\FormulaRefAuto{CanonicalProfileFamilies}[def];
\FormulaRefAuto{CanonicalProfileMapDefinition}[def]}
&
\textbf{Bijektion und Häufigkeitstransport}
\UebFormel{\begin{aligned}
&\Phi:A\bij\mathcal U,\qquad
 \UnionClosedFam(\mathcal U).
\end{aligned}}Ist \(A\ne\{0\}\), dann ist \(\mathcal U\) eine Frankl-Familie,
und\UebFormel{\begin{aligned}
&\HalfOcc{j}{\mathcal U}\ \vdash\
 \FranklSemi{A,\JoinOp,0}.
\end{aligned}}
\UebQuellen{\FormulaRefAuto{CanonicalProfileBijection}[thm];
\FormulaRefAuto{CanonicalProfileFamilyUnionClosed}[thm];
\FormulaRefAuto{CanonicalUnionFamilyIsFranklFamily}[thm];
\FormulaRefAuto{CanonicalHalfOccurrenceTransport}[thm]} \\
\addlinespace[.8ex]

\textbf{Komplementduale Konstruktion}
Für eine Frankl-Familie \(M\):\UebFormel{\begin{aligned}
&U=\bigcup M,\\
&L=\{U\setminus X\mid X\in M^\circ\},\\
&X\JoinOp Y=\bigcap\mathcal Z_{X,Y},\\
&\mathcal Z_{X,Y}=\{Z\in L\mid X\cup Y\subseteq Z\}.
\end{aligned}}
\UebQuellen{\FormulaRefAuto{FranklComplementConstruction}[def]}
&
\textbf{Ein endlicher Halbverband mit Null}
\UebFormel{\begin{aligned}
&\Finite L,\qquad\InterClosedFam L,\\
&\ZeroJoinSemi{L,\JoinOp,\varnothing},\\
&X,Y\in L\ \vdash\
 (X\le Y\leftrightarrow X\subseteq Y).
\end{aligned}}
\UebQuellen{\FormulaRefAuto{FranklComplementFamilyFinite}[thm];
\FormulaRefAuto{FranklComplementFamilyIntersectionClosed}[thm];
\FormulaRefAuto{FranklComplementFiniteLattice}[thm];
\FormulaRefAuto{FranklComplementOrderRecovered}[thm]} \\
\addlinespace[.8ex]

\textbf{Private Punkte}
In der vorigen Komplementkonstruktion besitzt jedes
\(\JoinOp\)-irreduzible \(P\in L\) einen privaten Punkt \(a\in P\):\UebFormel{\begin{aligned}
&Q\in L\ \vdash\
 (a\in Q\leftrightarrow P\le Q).
\end{aligned}}Der Hauptfilter \(H=\uparrow P\) ist damit genau die
Teilfamilie der \(a\) enthaltenden Mengen.
\UebQuellen{\FormulaRefAuto{JoinIrreduciblePrivatePoint}[thm]}
&
\textbf{Rücktransport zur Mengenform}
Für einen solchen privaten Punkt gilt:\UebFormel{\begin{aligned}
&\exists f:H\inj(L\setminus H)\\
&\qquad\vdash\HalfOcc{a}{M}.
\end{aligned}}Ein Irreduzibles mit höchstens halbgroßem Hauptfilter liefert
somit ein halbhäufiges Element der ursprünglichen Familie.
\UebQuellen{\FormulaRefAuto{FranklComplementPrivatePointTransport}[thm]} \\
\addlinespace[.8ex]

\textbf{Die beiden globalen Aussagen}
Die Mengenform lautet\UebFormel{\begin{aligned}
&\FranklSetStatement\ \leftrightarrow\\
&\quad\forall M\\
&\qquad(\FranklFam M\to\Frankl{M}).
\end{aligned}}Die Halbverbandsform \(\FranklJoinStatement\) fordert
\(\FranklSemi{A,\JoinOp,0}\) für jeden endlichen Halbverband mit Null
und mindestens einem von Null verschiedenen Element.
\UebQuellen{\FormulaRefAuto{FranklGlobalSetStatement}[def];
\FormulaRefAuto{FranklGlobalJoinStatement}[def]}
&
\textbf{Bewiesene Äquivalenz, offene Vermutung}
\UebFormel{\begin{aligned}
&\FranklSetStatement\leftrightarrow\FranklJoinStatement.
\end{aligned}}Der Satz beweist die beiden Transporte zwischen den globalen
Aussagen. Er beweist keine dieser Aussagen für sich.
\UebQuellen{\FormulaRefAuto{FranklSetImpliesSemilattice}[thm];
\FormulaRefAuto{FranklSemilatticeImpliesSet}[thm];
\FormulaRefAuto{FranklSetSemilatticeEquivalence}[thm]} \\
\addlinespace[.8ex]
\end{BandUebersicht}

